\subsection{Comparación con otras alternativas}

NFS convive con otras tecnologias de almacenamiento en red. SMB/CIFS cumple un rol similar en el mundo Windows: compartir archivos y directorios sobre la red. En ambientes Linux y UNIX, NFS suele integrarse de forma mas natural con permisos POSIX, UID/GID y herramientas del sistema. En organizaciones mixtas, ambos protocolos pueden coexistir sobre la misma infraestructura NAS.

iSCSI y Fibre Channel pertenecen al mundo del almacenamiento a bloques. Estas opciones son adecuadas cuando la aplicacion espera ver un disco remoto, cuando se requiere multipathing clasico o cuando se busca latencia mas deterministica. NFS, en cambio, entrega archivos ya administrados por el servidor. El almacenamiento de objetos, como S3, propone otro modelo: acceso mediante APIs, no mediante un sistema de archivos POSIX montado.

\subsection{NFS en la nube}

La computacion en la nube no elimino NFS; en muchos casos lo transformo en un servicio administrado. Amazon EFS, Google Filestore y Azure NetApp Files ofrecen almacenamiento de archivos accesible mediante NFS para instancias, contenedores o servicios gestionados. Aunque el proveedor abstrae parte de la alta disponibilidad y el mantenimiento, los conceptos de red siguen presentes: subredes, grupos de seguridad, latencia entre zonas, throughput provisionado y politicas de acceso.
